\begin{helper}[The subquadratic signed power lies in $W^{1,2}$ of every inner ball]
Let $n\in\mathbb{N}$ and let $p\in\R$ with $1<p<2$; in particular
$p/2-1<0$ and $p-2<0$.  All powers of nonnegative reals below are real powers,
so that $0^{p/2-1}=0$ for the negative exponent $p/2-1$.
$\operatorname{supp}$ denotes the closed support.  Let
$L_c,C_2,\varepsilon,\mu,\lambda\in\R$, let $g\colon\R^n\to\R^n$, let $Dg$
assign to each index $i$ a field $Dg^i\colon\R^n\to\R^n$, let $x_0\in\R^n$, let
$R\in\R$ and let $k$ be an index in $\{1,\dots,n\}$.  Assume
\begin{enumerate}
\item[(H1)] $0<\varepsilon$, $0<\mu$, $\lambda=\mu/2$, $0<\lambda$ and
  $0\le C_2$;
\item[(H2)] $|g^i|\le\mu$ almost everywhere on $B_R(x_0)$, for every index $i$;
\item[(H3)] (the radial cut-off) for all radii $0<\rho_0<\rho_1$ there is
  $\zeta\colon\R^n\to\R$ which is $C^\infty$ with compact support,
  $\operatorname{supp}\zeta\subseteq B_{\rho_1}(x_0)$, $0\le\zeta\le1$,
  $\zeta=1$ on $\overline B_{\rho_0}(x_0)$ and
  $\sum_i(D_i\zeta)^2\le L_c/(\rho_1-\rho_0)^2$ everywhere;
\item[(H4)] $\varepsilon+|g|^2\le(1+4n)\lambda^2$ almost everywhere on
  $B_R(x_0)$, and $\varepsilon\le\varepsilon+|g|^2$ everywhere;
\item[(H5)] each $g^i$ is almost everywhere strongly measurable on $B_R(x_0)$,
  and so is $Dg^k$.
\end{enumerate}
Let $\pi,G\colon\R\times\R\to\R$, written $\pi_\eta$ and $G_\eta$, satisfy
\begin{enumerate}
\item[(H6)] $G_\eta(s)=\int_0^s\frac p2(|t|+\eta)^{\frac p2-1}dt$ for all
  $\eta,s$;
\item[(H7)] $0<|s|+\eta$ for all $s$ and all $\eta>0$;
\item[(H8)] $G_\eta$ is $C^1$ for every $\eta>0$;
\item[(H9)] $|\pi_\eta(g^k)|\le\mu^{p-1}$ almost everywhere on $B_R(x_0)$, for
  every $\eta>0$;
\item[(H10)] for every $\eta>0$ and every $T$ with $0<T<R$,
  $v_\eta:=G_\eta(g^k)\in W^{1,2}(B_T(x_0))$ with weak gradient
  $V_\eta:=\frac p2(|g^k|+\eta)^{\frac p2-1}Dg^k$ (\noderef{MemW1pOn}).
\end{enumerate}
Let $e_0\in\R$ and $\eta_\bullet\colon\mathbb{N}\to\R$ satisfy
\begin{enumerate}
\item[(H11)] $0<e_0\le1$, $\eta_j=e_0/(j+1)$, $\eta_j>0$ and
  $\eta_j^2<\varepsilon$ for every $j$, and $\eta_j\to0$.
\end{enumerate}
Finally assume the Caccioppoli estimate at the regularisation level:
\begin{enumerate}
\item[(H12)] for every $\eta>0$ with $\eta^2<\varepsilon$, every $\ell\in\R$ and
  every $\zeta$ that is $C^\infty$ with compact support and
  $\operatorname{supp}\zeta\subseteq B_R(x_0)$: the function
  $\bigl(\varepsilon+|g|^2\bigr)^{\frac p2-1}\|V_\eta\|^2\zeta^2$ is integrable
  on $B_R(x_0)$, and, writing
  $A^{\pm}_\ell(\eta)$ for $B_R(x_0)\cap\{\pm(\pi_\eta(g^k)-\ell)>0\}$,
  \[
  \int_{A^{\pm}_\ell(\eta)}
  \bigl(\varepsilon+|g|^2\bigr)^{\frac p2-1}\|V_\eta\|^2\zeta^2\,dx
  \ \le\ C_2\int_{B_R(x_0)}
  \bigl((\pm(\pi_\eta(g^k)-\ell))^{+}\bigr)^2\sum_i(D_i\zeta)^2\,dx .
  \]
\end{enumerate}
Then for every $T$ with $0<T<R$,
\[
v:=|g^k|^{\frac p2-1}g^k\ \in\ W^{1,2}\bigl(B_T(x_0)\bigr)
\quad\text{with weak gradient}\quad
V:=\tfrac p2|g^k|^{\frac p2-1}Dg^k
\qquad(\noderef{MemW1pOn}).
\]
\end{helper}
\begin{proof}
Write $B_\rho=B_\rho(x_0)$ and fix $T$ with $0<T<R$.  From $1<p<2$ we get
$p>0$, $p/2-1<0$ and $p-2<0$.

\emph{Step 1: the cut-off and the lower bound on the weight.}  Put
$T'=(T+R)/2$, so $T<T'<R$, and let $\zeta$ be the cut-off of (H3) at
$\rho_0=T$, $\rho_1=T'$: it is $C^\infty$ with compact support,
$\operatorname{supp}\zeta\subseteq B_{T'}\subseteq B_R$, $0\le\zeta\le1$,
$\zeta=1$ on $\overline B_T$ and $\sum_i(D_i\zeta)^2\le L_c/(T'-T)^2$.  By (H4)
and $p/2-1<0$, the map $t\mapsto t^{p/2-1}$ being decreasing on $(0,\infty)$,
\[
\bigl(\varepsilon+|g|^2\bigr)^{\frac p2-1}
\ \ge\ W_\ast:=\bigl((1+4n)\lambda^2\bigr)^{\frac p2-1}>0
\qquad\text{almost everywhere on }B_R ,
\tag{1}
\]
the base being positive by (H1).

\emph{Step 2: a bound on $\int_{B_T}\|V_\eta\|^2$ uniform in $\eta$.}  Take
$\ell_0=-\mu^{p-1}-1$ in (H12), with the $+$ sign and this $\zeta$.  By (H9),
$\pi_\eta(g^k)\ge-\mu^{p-1}>\ell_0$ almost everywhere on $B_R$, so
$A^{+}_{\ell_0}(\eta)$ is all of $B_R$ up to a null set and \emph{no splitting
of $B_R$ by the sign of $g^k$ is needed}; moreover
$\bigl((\pi_\eta(g^k)-\ell_0)^{+}\bigr)^2\le(\mu^{p-1}+|\ell_0|)^2$ almost
everywhere there.  Hence the right-hand side of (H12) is at most
\[
\mathcal K:=C_2\,(\mu^{p-1}+|\ell_0|)^2\,\frac{L_c}{(T'-T)^2}\,|B_R|<\infty ,
\]
using $\sum_i(D_i\zeta)^2\le L_c/(T'-T)^2$ and $|B_R|<\infty$.  On the left,
$\zeta=1$ on $\overline B_T\supseteq B_T$ and the integrand is nonnegative, so
by (1) and monotonicity of the integral (Preamble)
\[
W_\ast\int_{B_T}\|V_\eta\|^2dx
\ \le\ \int_{B_R}\bigl(\varepsilon+|g|^2\bigr)^{\frac p2-1}\|V_\eta\|^2\zeta^2dx
\ \le\ \mathcal K ,
\qquad\text{so}\qquad
\int_{B_T}\|V_\eta\|^2dx\ \le\ \frac{\mathcal K}{W_\ast}
\tag{2}
\]
for every $\eta>0$ with $\eta^2<\varepsilon$, the integrability of the middle
integrand being the first clause of (H12).

\emph{Step 3: one monotone limit, giving both $V\in L^2(B_T)$ and
$Dg^k=0$ almost everywhere on $\{g^k=0\}$.}  Since $p/2-1<0$, the map
$\eta\mapsto(|s|+\eta)^{p/2-1}$ is nonincreasing in $\eta$ for each fixed $s$,
so $\|V_{\eta_j}\|^2$ is nondecreasing in $j$ by (H11), $\eta_j$ being
nonincreasing.  Put $S(x)=\sup_j\|V_{\eta_j}(x)\|^2\in[0,\infty]$.  By the
monotone convergence theorem (Preamble) and (2) at $\eta=\eta_j$, which (H11)
licenses,
\[
\int_{B_T}S\,dx=\lim_{j\to\infty}\int_{B_T}\|V_{\eta_j}\|^2dx
\ \le\ \frac{\mathcal K}{W_\ast}<\infty ,
\]
so $S<\infty$ almost everywhere on $B_T$.  \emph{That single finiteness
statement does both jobs.}  First, at a point where $g^k=0$ and $Dg^k\ne0$ one
has $\|V_{\eta_j}\|^2=\frac{p^2}4(\eta_j)^{p-2}\|Dg^k\|^2$, and
$(\eta_j)^{p-2}\to+\infty$ because $\eta_j\downarrow0$ and $p-2<0$; so
$S=+\infty$ there, and therefore the set of such points is null:
\[
Dg^k=0\quad\text{almost everywhere on }\{g^k=0\}.
\tag{3}
\]
Second, $\|V_{\eta_j}\|^2\uparrow\|V\|^2$ almost everywhere on $B_T$: on
$\{g^k\ne0\}$ because $(|g^k|+\eta_j)^{p/2-1}\uparrow|g^k|^{p/2-1}$, and on
$\{g^k=0\}$ because both sides vanish almost everywhere there, $V=0$ since
$0^{p/2-1}=0$ for the negative exponent and $V_{\eta_j}=0$ almost everywhere by
(3).  Hence $\|V\|^2\le S$ almost everywhere and
$\int_{B_T}\|V\|^2\le\mathcal K/W_\ast<\infty$, so $V\in L^2(B_T)$;
measurability of $V$ follows from (H5) with $B_T\subseteq B_R$.

\emph{Step 4: $v\in L^2(B_T)$.}  By (H2), $|g^k|\le\mu$ almost everywhere on
$B_R$, and $|v|=\bigl||g^k|^{p/2-1}g^k\bigr|\le|g^k|^{p/2}\le\mu^{p/2}$ almost
everywhere there, the first inequality holding at every real $t$ in the form
$\bigl||t|^{p/2-1}t\bigr|\le|t|^{p/2}$ --- with equality off $0$ and both sides
$0$ at $t=0$.  Since $B_T$ has finite measure, a bounded almost everywhere
measurable function lies in $L^2(B_T)$, and $v$ is measurable by (H5).

\emph{Step 5: the closed form of $G_\eta$ and the pointwise limits.}  By (H6)
and (H7) the integrand $t\mapsto\frac p2(|t|+\eta)^{p/2-1}$ is continuous, so
for $s\ge0$ one has $|t|=t$ on $[0,s]$ and the fundamental theorem of calculus
in interval-integral form (Preamble), with antiderivative $(t+\eta)^{p/2}$,
gives $G_\eta(s)=(s+\eta)^{p/2}-\eta^{p/2}$; symmetrically
$G_\eta(s)=-\bigl((|s|+\eta)^{p/2}-\eta^{p/2}\bigr)$ for $s\le0$.  \emph{The
closed form is what makes the limit elementary}: as $j\to\infty$,
$(|s|+\eta_j)^{p/2}\to|s|^{p/2}$ by continuity and $\eta_j^{p/2}\to0$ because
$p/2>0$, so
\[
G_{\eta_j}(s)\ \longrightarrow\ |s|^{p/2}\operatorname{sgn}s=|s|^{\frac p2-1}s
\qquad(s\in\R),
\]
that is, $v_{\eta_j}\to v$ pointwise.  Also
$|G_\eta(s)|\le(|s|+\eta)^{p/2}$, so with (H2) and $\eta_j\le e_0\le1$ from
(H11), $|v_{\eta_j}|\le(\mu+1)^{p/2}$ almost everywhere on $B_R$.  And
$V_{\eta_j}\to V$ pointwise almost everywhere on $B_T$ componentwise, by the
computation of Step 3, with the domination
$|(V_{\eta_j})_i|\le\|V_{\eta_j}\|\le\|V\|$ almost everywhere there.

\emph{Step 6: dominated convergence in the identity.}  Let $\varphi$ be
$C^\infty$ with compact support and $\operatorname{supp}\varphi\subseteq B_T$,
and fix an index $i$.  By (H10) at $\eta=\eta_j$ and the radius $T$,
\[
\int_{B_T}v_{\eta_j}\,D_i\varphi\,dx=-\int_{B_T}(V_{\eta_j})_i\,\varphi\,dx
\qquad(j\in\mathbb{N}),
\]
this being the weak-gradient identity of (H10) read at the test function
$\varphi$ (\noderef{HasWeakGradientOn}).  Let $j\to\infty$.  On the left,
$v_{\eta_j}\to v$ pointwise by Step 5 and
$|v_{\eta_j}D_i\varphi|\le(\mu+1)^{p/2}|D_i\varphi|$ almost everywhere, a
bounded function on the finite-measure set $B_T$ since $D_i\varphi$ is
continuous with compact support, hence integrable.  On the right,
$(V_{\eta_j})_i\to V_i$ almost everywhere and
$|(V_{\eta_j})_i\varphi|\le\|V\|\,|\varphi|$ almost everywhere by Step 5, which
is integrable on $B_T$ because $\|V\|\in L^2(B_T)$ by Step 3, $\varphi$ is
bounded and $|B_T|<\infty$, so H\"older on a set of finite measure (Preamble)
applies.  Dominated convergence (Preamble) therefore passes to the limit on
both sides and gives
\[
\int_{B_T}v\,D_i\varphi\,dx=-\int_{B_T}V_i\,\varphi\,dx .
\]
Since $\varphi$ and $i$ were arbitrary, $V$ is the weak gradient of $v$ on
$B_T$ (\noderef{HasWeakGradientOn}).  Together with $v\in L^2(B_T)$ from
Step 4 and $V\in L^2(B_T)$ from Step 3, that is $v\in W^{1,2}(B_T)$ with weak
gradient $V$ (\noderef{MemW1pOn}), which is the assertion.
\end{proof}
